\section*{Author Section 5: fully-derandomized switching lemma}
Before we show the proof of Lemma 9, we state and prove the simpler version of the classical switching lemma in this section. The following statement is the full version of Lemma 2 in Introduction. Lemma 2 follows from Lemma 11 trivially.
\paragraph{Lemma 11.} Let $k,m\ge1$ be integers, and $\epsilon,p>0$ be reals. Let $F=\bigvee_{i=1}^mC_i$ be a size-$m$ $k$-DNF over inputs $\{0,1\}^n$. For all $t\ge1$, let $(\boldsymbol\Lambda,\boldsymbol x)$ be a joint random variable such that:
\begin{itemize}\item $\boldsymbol\Lambda$ is a $(t+k)$-wise $p$-bounded subset of $[n]$.
\item Conditioning on $\boldsymbol\Lambda$, $\boldsymbol x$ is a random string that $\epsilon$-fools CNF of size at most $m$.\end{itemize}
Consider the random restriction $F|_{\Lambda[x,\star]}$, we have:
\[\Pr_{\Lambda,x}[\mathrm{DT}(F|_{\Lambda[x,\star]})>t]\le(10kp)^t+(4m)^{t+k}\epsilon.\]
Understanding the proof of Lemma 11 is necessary to read the proof of Lemma 9. On the other hand, once Lemma 11 is established, we can prove Lemma 9 using a rather similar strategy. The rest of the section is devoted to the proof of Lemma 11. We will first introduce the important concept of ``canonical decision tree'' ([14,27,25,26]) in Section 5.1. We prove for the case that $x$ is truly random (and only $\Lambda$ is pseudorandom) in Section 5.2, and then argue how to prove for pseudorandom $x$ in Section 5.3. Throughout the whole section, we always use $F=\bigvee_{i=1}^mC_i$ to denote the $k$-DNF that we are analyzing. For every $i\in[m]$ let $V_i\subseteq[n]$ be the variables involved in the term $C_i$.
\subsection*{5.1 Canonical decision tree}
For a DNF $F=\bigvee_{i=1}^mC_i$, denote by $T_F$ the canonical decision tree of $F$, whose construction is shown in Algorithm 1. We use $\mathrm{CDT}(F)$ to denote the depth of the canonical decision tree for $F$. Since canonical decision tree is one particular decision tree of $F$, its depth must be no less than $\mathrm{DT}(F)$. To prove Lemma 11, we will analyze canonical decision trees, and show that with high probability over the random restriction $\rho$, we have $\mathrm{CDT}(F|_\rho)<t$.
\paragraph{Algorithm 1: Canonical Decision Tree.}
Input: A DNF $F=\bigvee_{i=1}^mC_i$, black-box access to a string $\alpha\in\{0,1\}^n$.
\begin{quote}\textbf{initialize:} $j^*\leftarrow0$; $x\leftarrow(\star)^n$.\\
\textbf{while} $j^*<m$ \textbf{do}\\
\hspace*{1em}Find the first $j>j^*$ such that $C_j(x)\not\equiv0$. If no such $j$ exists, exit the loop.\\
\hspace*{1em}$B_j\leftarrow$ the set of unknown variables in $C_j$.\\
\hspace*{1em}Query $\alpha_{B_j}$.\\
\hspace*{1em}Set $x_{B_j}\leftarrow\alpha_{B_j}$.\\
\hspace*{1em}\textbf{if} $C_j(x)=1$ \textbf{then return} $1$.\\
\hspace*{1em}$j^*\leftarrow j$.\\
\textbf{return} $0$.\end{quote}
\subsection*{5.2 Proof when $x$ is truly random}
\paragraph{Defining witness.} Let $\rho=\Lambda[x,\star]$ being a bad restriction, under which $\mathrm{CDT}(F|_\rho)\ge t$. Consider simulating the canonical decision tree $T_{F|_\rho}$. We know that on some inputs $\alpha\in\{0,1\}^n$, $T_{F|_\rho}$ fails to output the decision after making $(t-1)$ queries. We choose one such $\alpha$ and simulate $T_{F|_\rho}$ until it makes at least $t$ queries. The ``running transcript'' of $T_{F|_\rho}$ on $\alpha$ is naturally a witness to the fact that $\mathrm{CDT}(F_\rho)\ge t$. We formalize this idea in the following definition.
\paragraph{Definition 12.} Let $F=\bigvee_{i=1}^mC_i$ be the $k$-DNF and $\rho\in\{0,1,\star\}^n$ be a restriction. Let $t\ge1$. Consider a tuple $(r,\ell_i,s_i,B_i,\alpha_i)$, where:
\begin{enumerate}\item $r\in[1,t]$ is an integer.
\item $(\ell_1,\ldots,\ell_r)\in[m]^r$ is a list of increasing indices.
\item $(s_1,\ldots,s_r)$ is a list of positive integers such that $s:=\sum_{i=1}^rs_i\in[t,t+k-1]$.
\item $(B_1,\ldots,B_r)$ is a list of subsets of $[k]$. Moreover, for every $i\in[r]$, $|B_i|=s_i$.
\item $(\alpha_1,\ldots,\alpha_r)$ is a list of binary strings. For every $i\in[r]$, $|\alpha_i|=s_i$.\end{enumerate}
We call $(r,\ell_i,s_i,B_i,\alpha_i)$ a $t$-witness for $\rho$, if there exists $\alpha\in\{0,1\}^n$ such that:
\begin{itemize}\item When we run $T_{F|_\rho}$ on $\alpha$, for every $i\in[r]$, $C_{\ell_i}$ is the $i$-th term queried by $T_{F|_\rho}$.
\item By the time $T_{F|_\rho}$ issues the $i$-th set of query, exactly $s_i$ variables in $C_{\ell_i}$ are not known, and their ``relative positions'' in $V_{\ell_i}$ are specified by $B_i$.
\item The response to the $i$-th set of query is $\alpha_i$.\end{itemize}
We define the size of the witness $(r,\ell_i,s_i,B_i,\alpha_i)$ as $s:=\sum_{i=1}^rs_i$.
\paragraph{Notation convention.} When $B_i$ and $\alpha_i$ are associated with a term $C_{\ell_i}$, sometimes we will slightly abuse notation by using $B_i$ to refer to the set of variables it corresponds to in $V_{\ell_i}$ (Recall $V_{\ell_i}$ is the set of variables that appear in $C_{\ell_i}$), and use $\alpha_i$ to denote a partial assignment $\alpha_i\in\{0,1,\star\}^n$ in which we only assign the $B_i$ part of $\alpha_i$ and leave other coordinates unfixed. In this section, we will mostly use the term ``witness'' to denote $t$-witness, since we are always analyzing for a fixed $t$.
It is easy to see that if $\mathrm{DT}(F|_\rho)\ge t$, there must be a witness for $\mathrm{CDT}(F|_\rho)\ge t$.\footnote{However, we note that the converse may not be true, since our witness can only refute the existence of shallow canonical decision trees.} By now, a natural idea to bound the probability of picking a bad restriction would be (1) enumerating every possible witness, (2) calculating the probability of a random $\rho$ having such a witness, and (3) union-bounding over them. Unfortunately this is too expensive for us: we have at least $\binom mt$ choices of the list $(\ell_i)$. In order for this approach to be meaningful, we have to bound the probability that a random restriction has a particular witness by $m^{-t}$, which seems very hard, if not impossible. It turns out we can avoid the enumeration of $(\ell_i)$ part in the witness. To succinctly describe the idea, let us define partial witnesses first.
\paragraph{Definition 13.} Let $F=\bigvee_{i=1}^mC_i$ be the $k$-DNF and $\rho\in\{0,1,\star\}^n$ be a restriction. Let $t\ge1$. Consider a tuple $(r,s_i,B_i,\alpha_i)$. We call $(r,s_i,B_i,\alpha_i)$ a partial $t$-witness for $\rho$, if there exists $(\ell_1,\ldots,\ell_r)$ such that $(r,\ell_i,s_i,B_i,\alpha_i)$ is a $t$-witness for $\rho$.
\paragraph{Remark 14.} Here we make an important observation: if $(r,s_i,B_i,\alpha_i)$ is a partial witness for $\rho$, then there is only one valid list $(\ell_i)$ which makes $(r,\ell_i,s_i,B_i,\alpha_i)$ a witness for $\rho$ (To see this, first note that $\ell_1$ is fixed. After querying $\ell_1$ and getting the response $\alpha_1$, use induction). This observation will prove useful in Section 5.3 and Section 6.
Then, the proof goes by enumerating $r,s_i,B_i,\alpha_i$ and bounding the following:
\[\Pr_\rho[\mathrm{DT}(F|_\rho)\ge t]\le\sum_{(r,s_i,B_i,\alpha_i)}\Pr_\rho[(r,s_i,B_i,\alpha_i)\text{ is a partial witness for }\rho].\tag{5}\]
Fixing an $s\in[t,t+k-1]$, there are at most $(4k)^s$ possible partial witnesses $(r,s_i,B_i,\alpha_i)$ of size $s$. For each of them, if we can bound the probability by, say, $(2p)^s$, then the lemma is proved.
\paragraph{The witness searcher.} Now, given a restriction $\rho$ and a candidate partial witness $(r,s_i,B_i,\alpha_i)$, how can we decide if there is a list $(\ell_i)_i$ such that $(r,\ell_i,s_i,B_i,\alpha_i)$ constitutes a witness for $\rho$? We will do so by designing and using a simple procedure, which we call the witness searcher. More specifically, our searcher receives as input a restriction $\rho$, a partial witness $(r,s_i,B_i,\alpha_i)$ and an advice string $y\in\{0,1\}^n$. It outputs either a complete witness $(r,\ell_i,s_i,B_i,\alpha_i)$, or a special ERROR symbol. Its description in shown in Algorithm 2. In the following, we use $S$ to denote Algorithm 2.
\paragraph{Algorithm 2: Witness Searcher.} Input: A DNF $F=\bigvee_{i=1}^mC_i$, a restriction $\rho\in\{0,1,\star\}^n$, a partial witness $(r,s_i,B_i,\alpha_i)$, an advice $y\in\{0,1\}^n$.
\begin{quote}\textbf{initialize:} $z\leftarrow\rho\circ y$; $j^*\leftarrow0$; $c\leftarrow1$.\\
\textbf{while} $c\le r$ \textbf{do}\\
\hspace*{1em}Find the first $j>j^*$ such that $C_j$ is satisfied by $z$. If no such $i$ exists, return ERROR.\\
\hspace*{1em}Set $\ell_c\leftarrow j$, and associate $B_c,\alpha_c$ with $C_{\ell_c}$.\\
\hspace*{1em}Replace the $B_c$ part of $z$ with $\alpha_c$. That is, $z\leftarrow\rho\circ\alpha_1\circ\cdots\circ\alpha_c\circ y$.\\
\hspace*{1em}$c\leftarrow c+1$; $j^*\leftarrow j$.\\
\textbf{return} $(r,\ell_i,s_i,B_i,\alpha_i)$.\end{quote}
\paragraph{Lemma 15.} If $(r,s_i,b_i,\alpha_i)$ is a partial witness for $\rho$, then there exists an advice $y$ that makes $S$ find $(\ell_i)_{i=1}^r$. More importantly, on a uniformly random $y\sim U_n$, $S$ finds $(\ell_1,\ldots,\ell_r)$ with probability exactly $2^{-s}$.
\begin{proof}If $(r,s_i,\ell_i,b_i,\alpha_i)$ is a witness for $\rho$, it records a running transcript of $T_{F|_\rho}$ on some input $\alpha\in\{0,1\}^n$. This mean that $C_{\ell_1}$ is the first term queried by $T_{F|_\rho}$, which implies:
\begin{enumerate}\item Every term before $C_{\ell_1}$ is falsified by $\rho$ (because $T_{F|_\rho}$ skipped them).
\item $C_{\ell_1}$ is consistent with $\rho$ (because $T_{F|_\rho}$ was uncertain of the value of $C_{\ell_i}(\alpha)$.\end{enumerate}
Now, consider running $S$ on $(r,s_i,b_i,\alpha_i)$, $S$ will also skip all the terms before $C_{\ell_1}$. At $C_{\ell_1}$, since $\rho$ is consistent with $C_{\ell_1}$, the random string $\rho\circ y$ satisfies $C_{\ell_1}$ with probability $2^{-|B_1|}$. Conditioning on this happened, in the later execution, we modify $z$ and replace the part corresponding to $B_1$ with $\alpha_1$.
At this point, we go back and inspect the execution of $T_{F|_\rho}$ on $\alpha$. Since $(r,s_i,\ell_i,b_i,\alpha_i)$ is the running transcript of $T_{F|_\rho}$ on $\alpha$, we know that $\alpha_1$ was the response that $T_{F|_\rho}$ received from querying alive variables in $C_{\ell_1}$. Since $T_{F|_\rho}$ issued its second bunch of queries to alive variables in $C_{\ell_2}$, it implies that terms from $C_{\ell_1+1}$ to $C_{\ell_2-1}$ are all falsified by $\rho\circ\alpha_1$, and $C_{\ell_2}$ is consistent with $\rho\circ\alpha_1$. Then, it follows that $\rho\circ\alpha_1\circ y$ satisfies $C_{\ell_2}$ with probability $2^{-|B_2|}$. Conditioning on this happened, the searcher will proceed with string $\rho\circ\alpha_1\circ\alpha_2\circ y$ and we can again consider the execution of $T_{F|_\rho}$ after querying $C_{\ell_2}$. We do this argument so on and so forth, until we have identified all of $r$ indices $\ell_1,\ldots,\ell_r$ and exit the procedure. In summary, we have shown there exists $y$ which makes $S$ find the list $(\ell_i)_{i=1}^r$, and the probability of sampling such a $y$ is $2^{-\sum_{i=1}^c|B_i|}=2^{-s}$ as desired.\end{proof}
\paragraph{Decoupling.} If we inspect the execution of $S$ carefully, we can notice that it only needs to know the string $z=\rho\circ y$ to work. In particular, it does not need to know which part of $z$ is fixed in the restriction $\rho$. Therefore, we can revise $S$ to get a searcher $S'$: the input to $S'$ is now a string $z$ and $(r,s_i,B_i,\alpha_i)$. Otherwise it runs identically the same way as $S$. We denote the output of $S'$ as $S'(z,(r,s_i,B_i,\alpha_i))$. By Lemma 15, we have
\begin{multline*}\Pr_{\rho\sim R}[(r,s_i,B_i,\alpha_i)\text{ is a partial witness for }\rho]\\\le2^s\Pr_{\rho\sim R,y\sim U_n}[S'(\rho\circ y,(r,s_i,B_i,\alpha_i))\text{ is a witness for }\rho].\end{multline*}
Denote $\rho=\Lambda[x,\star]$. We observe that
\[\mathbf1_{\{(r,\ell_i,s_i,B_i,\alpha_i)\text{ is a witness for }\rho\}}\le\mathbf1_{\{(\bigcup_{j=1}^rB_j)\subseteq\Lambda\}}.\tag{6}\]
For a $(t+k)$-wise $p$-bounded set $\Lambda$, the event on the right hand side holds with probability $p^{\sum_j|B_j|}$. Then, we have
\begin{align*}&\Pr_{\rho\sim\Lambda[x,0],y\sim U_n}[S'(\rho\circ y,(r,s_i,B_i,\alpha_i))\text{ is a witness for }\rho]\\&=\mathbb E_\Lambda\mathbb E_{x,y\sim U_n}[\mathbf1_{\{S'(\Lambda[x,y],(r,s_i,B_i,\alpha_i))\text{ is a witness for }\rho\}}]\\&=\mathbb E_{z\sim U_n}\left[\Pr_\Lambda[S'(z,(r,s_i,B_i,\alpha_i))\text{ is a witness for }\rho]\right]\\&\le p^s,\end{align*}
where the third equality is due to that $\Lambda[x,y]$ is distributed as $U_n$ when $x,y\sim U_n$, and the last inequality holds by (6) and the $(t+k)$-wise $p$-bounded property of $\Lambda$.
\paragraph{Wrapping-up.} We finish the proof by enumerating all $(r,s_i,B_i,\alpha_i)$ and taking a summation.
\begin{align*}&\Pr_{\rho\sim\Lambda[x,\star]}[\mathrm{DT}(F|_\rho)\ge t]\\&\le\sum_{(r,s_i,B_i,\alpha_i)}\Pr_{\rho\sim\Lambda[x,\star]}[(r,s_i,B_i,\alpha_i)\text{ is a partial witness for }\rho]\\&\le\sum_{(r,s_i,B_i,\alpha_i)}2^s\mathbb E_{z\in U_n}\left[\Pr_\Lambda[S'(z,(r,s_i,B_i,\alpha_i))\text{ is a witness for }\Lambda[x,\star]]\right]\\&\le\sum_{s=t}^{t+k}2^s(4k)^sp^s\\&\le(10kp)^t.\tag{7}\end{align*}
\subsection*{5.3 Proof when $x$ is pseudorandom}
Now we consider the case that $x$ is not truly random. There could even be correlation between $x$ and $\Lambda$. Our only requirement for $x$ is that, for every fixed $\Lambda$, $x$ is a pseudorandom string that $\epsilon$-fools CNFs of size at most $m$. First of all, by Remark 14, the following equation is established.
\begin{multline*}\mathbf1_{\{(r,s_i,B_i,\alpha_i)\text{ is a partial witness for }\rho\}}\\=\sum_{\ell_i}\mathbf1_{\{(r,\ell_i,s_i,B_i,\alpha_i)\text{ is a witness for }\rho\}}.\tag{8}\end{multline*}
If we inspect the deduction in (7), it actually shows the following:
\[\sum_{(r,s_i,B_i,\alpha_i)}\sum_{\ell_i}\Pr_{\rho\sim\Lambda[U_n,\star]}[(r,\ell_i,s_i,B_i,\alpha_i)\text{ is a witness for }\rho]\le(10kp)^t.\tag{9}\]
Next, fixing the set $\Lambda$ and a tuple of $(r,\ell_i,s_i,B_i,\alpha_i)$, we consider the predicate
\[h_\Lambda^{(r,\ell_i,s_i,B_i,\alpha_i)}(x):=\mathbf1_{\{(r,\ell_i,s_i,B_i,\alpha_i)\text{ is a witness for }\Lambda[x,\star]\}}.\]
We omit the superscript and subscript of $h$ when they are clear from context. We claim that $h$ can be implemented by a CNF of size at most $m$. To see this, note that $h(x)=1$ if all of the following hold.
\begin{itemize}\item For every $j<\ell_1$, $C_j$ is falsified by $\Lambda[x,\star]$. For this to be true, there is at least one variable in $C_j$ that is not chosen by $\Lambda$ and is assigned to the opposite value to its appearance in $C_j$. Note that we can use an OR over relevant variables or their negations to verify if $C_j$ is falsified.
\item For $j=\ell_1$, $C_{\ell_1}$ cannot be falsified by $\Lambda[x,\star]$, which means all variables in $C_{\ell_1}$ that are not chosen by $\Lambda$ are assigned so that it does not contradict $C_{\ell_1}$. We can use an AND over those variables to verify this.
\item For every $\ell_1<j<\ell_2$, $C_j$ is falsified by $\Lambda[x,\alpha_1]$, which means that either $C_j$ is falsified by $\alpha_1$, or there is at least one variable in $C_j$ that is not chosen by $\Lambda$ and is assigned to the opposite value. If $C_j$ is falsified by $\alpha_1$, there is nothing to verify. Otherwise, we can verify it by an OR over relevant variables.
\item The same reasoning applies to $C_j$ for every $j\ge\ell_2$. Finally, after we finish the verification for $C_{\ell_r}$, we are done.\end{itemize}
In summary, we can implement $h$ by a CNF (AND of ORs). It is obvious that the size of the CNF is no larger than the size of $F$, which is $m$. Therefore, we conclude that $x$ $\epsilon$-fools $h$. That is,
\[\Pr_{x\sim U_n}[h_\Lambda^{(r,\ell_i,s_i,B_i,\alpha_i)}(x)]-\Pr_{x\sim\boldsymbol x}[h_\Lambda^{(r,\ell_i,s_i,B_i,\alpha_i)}(x)]\le\epsilon.\tag{10}\]
Finally, we have
\begin{align*}&\Pr_{\rho\sim\Lambda[\boldsymbol x,\star]}[\mathrm{DT}(F|_\rho)\ge t]\\&\le\sum_{(r,s_i,B_i,\alpha_i)}\Pr_{\rho\sim\Lambda[\boldsymbol x,\star]}[(r,s_i,B_i,\alpha_i)\text{ is a partial witness for }\rho]\\&\le\sum_{(r,s_i,B_i,\alpha_i)}\sum_{\ell_i}\Pr_{\rho\sim\Lambda[\boldsymbol x,\star]}[(r,\ell_i,s_i,B_i,\alpha_i)\text{ is a witness for }\rho]\\&\le\sum_{(r,s_i,B_i,\alpha_i)}\sum_{\ell_i}\epsilon+\Pr_{\rho\sim\Lambda[U_n,\star]}[(r,\ell_i,s_i,B_i,\alpha_i)\text{ is a witness for }\rho]\\&\le(10kp)^t+(4m)^{t+k}\epsilon.\end{align*}
Here, the first line is due to the argument in Section 5.2. The second line holds by (8). The third line is due to (10). The last line holds because (9) and there are at most $(4m)^{k+t}$ possible tuples $(r,\ell_i,s_i,B_i,\alpha_i)$.
